\subsection{The output of the application, panel by panel}
The flow of water and the uptake by roots described in Chapter 2 are shown by the layer \emph{Root structure} as pictures of the real output of the Daisy ensemble: the roots on six dates of a season, the depth at which half of the uptake occurs, and the rooting depth of every soil.

The upper panel is the density of roots in the profile through the year (darker means more roots) and the lower panel the water taken by the roots (darker green is more uptake, blue is irrigation). The red line is the date chosen, here 16 April, when the roots reach 30 cm.

\appfig{r1_0}{0.96}{Roots through the season: real model output, day by day, for the coarse loam in the weather of 1987.}

The six drawings are the root system on six dates, and the curve below is the rooting depth in each of the 20 years of weather. The roots grow from 23 cm in January to 94 cm in July and are cut back at harvest, which is the pattern assumed by the root term $S(z)$ of the flow equation.

\appfig{r1_1}{0.96}{Root structure: where the roots are and where the water is taken, with the rooting depth of every year of the simulation.}

For the day chosen the figure gives the root density, the water taken that day (green) and the water taken over the year (blue), with the depth of half and of 90 per cent of the uptake marked (here 17.5 and 55 cm); the water is taken from the upper roots where the soil is wet.

\appfig{r1_2}{0.96}{The day on the clock: the root density, the water taken today and the water taken in the year, at depth.}

The table gives the deepest roots and the depths of 50 and 90 per cent of the uptake for the ten soils, and the drawings show the roots of each soil on one day. Rooting depth differs by 11 cm among the soils (85 to 96 cm), and the model stops roots at 1 m.

\appfig{r1_3}{0.96}{Rooting depth in all ten soils, under conventional, reduced and no tillage, for the crop of the explorer.}

The graphs show the yield, irrigation, drainage and other results against the rooting depth for 1\,000 runs, and the table gives the result of one run (roots of 120 cm, 6 mm per application). Deeper roots raise the yield slowly and reduce the nitrogen leached, but the gain vanishes below about 1.2 m.

\appfig{r1_4}{0.96}{What your rooting-depth runs show: 1\,000 simulations on the coarse loam, with the rooting depth from 30 to 180 cm.}

The coloured band is the pressure head at the sensor depth through the year for the 20 simulated years: it falls from field capacity in spring to the dry end of the range in summer. The vertical drawings at the right show the profile of the day.

\appfig{v2_0}{0.96}{Soil pressure head through the year, January to December, every year 2019 to 2038.}

